\documentclass[10pt,a4paper]{amsart}
\usepackage[margin=19mm]{geometry}
\usepackage{amsmath,amssymb,mathtools}
\usepackage[scr=rsfs,cal=euler]{mathalfa}
\usepackage{xfrac}
\usepackage[unicode,hidelinks,bookmarks=false]{hyperref}
\newcommand{\ie}{i.e.\ }
\newcommand{\cf}{cf.\ }
\newcommand{\resp}{respectively\ }
\newcommand{\Subsection}{\subsection}
\providecommand{\nobreakdash}{\nobreak\hbox{-}}
\setlength{\emergencystretch}{2em}
\multlinegap0pt
\usepackage{bm}
\usepackage{bbm}
\usepackage{dsfont}
\usepackage{xspace}
\usepackage{cancel}
\usepackage{mathtools}
\usepackage{amsthm}
\makeatletter
\@ifundefined{theorem}{\newtheorem{theorem}{Theorem}}{}
\@ifundefined{lemma}{\newtheorem{lemma}[theorem]{Lemma}}{}
\makeatother
\DeclareMathOperator{\cone}{cone}
\newcommand{\GG}{G_{\overline{\chi}}}
\newcommand{\KK}{\mathbb{K}}
\newcommand{\ZZ}{\mathbb{Z}}
\newcommand{\ol}[1]{\overline{#1}}
\title[Root monoids and active algebraic groups]{Root monoids and active algebraic groups}
\author{Ekaterina Nistiuk (Presnova), Yulia Zaitseva}
\date{Source: arXiv:2510.16638v1 (CC BY 4.0)}
\begin{document}
\maketitle

\noindent\emph{Source and licence.} Root monoids and active algebraic groups. Version arXiv:2510.16638v1 (CC BY 4.0), distributed under the Creative Commons Attribution 4.0 International licence, as recorded in the arXiv licence metadata for this exact version and shown on the abstract page https://arxiv.org/abs/2510.16638v1. Full rights notice: licenses/arxiv-2510.16638-CC-BY-4.0.txt.\par\smallskip

\noindent\textit{Editorial scope.} Counted unit: the Theorem whose source label is \texttt{idem\_thm\_1} in the article (source0.tex lines 487--494, source label \texttt{idem\_thm\_1}), delivered together with the article's own complete original proof of it (source0.tex lines 496--557, one \texttt{proof} environment of 62 lines). No mathematics is written, repaired or re-derived: statement and proof are the article's own text line for line. Required context retained uncounted: tor:fun\_orb2 (lines 173--178); thm:mon (lines 279--285); thm:finalcom (lines 281--283). Every definition, display and label of those lines is retained unchanged. Omitted from this package and not redistributed: 1--172, 179--278, 284--486, 495--495, 558--874. Nothing inside a retained excerpt is omitted or altered: every retained body is delivered exactly as the article's lines are written, so no omission receipt of any kind accompanies this package. Citations: the retained excerpts carry no citation command at all, so no bibliography entry of the article is transcribed and no reference is redistributed. Reference adaptation: none was needed and none was made; every reference inside the delivered unit resolves inside it, so no unresolved reference or citation is produced. Printed slips kept literal: no printed word, symbol, hypothesis or number of the counted statement or of its proof was altered. Layout and class: the article's own class is \texttt{amsart} with packages including inputenc, fontenc, babel, amssymb, amscd, xy, geometry, dsfont, todonotes, mathrsfs, mathtools, graphicx, tikz, hyperref; the wrapper is the lane's pinned \texttt{amsart} editorial wrapper at 10pt on A4, which restates the theorem-like environments the retained text needs and adds the article's own macro definitions that the retained text needs (\texttt{\textbackslash GG}, \texttt{\textbackslash KK}, \texttt{\textbackslash ZZ}, \texttt{\textbackslash cone}, \texttt{\textbackslash ol}), with extra wrapper packages (bm, bbm, dsfont, xspace, cancel, mathtools). The delivered numbering is the unit's own. Assets: no retained span contains a figure environment, a graphics-inclusion command, a PDF-page inclusion, a table or a TikZ picture; no image file is copied into this package, the package declares no asset dependency and no image file is redistributed with it. Licence: version arXiv:2510.16638v1 of this article is distributed under the Creative Commons Attribution 4.0 International licence, as recorded in the arXiv licence metadata for this exact version and shown on the abstract page https://arxiv.org/abs/2510.16638v1. A full rights notice is delivered at licenses/arxiv-2510.16638-CC-BY-4.0.txt. Characters: every retained excerpt is pure ASCII, so the delivered engine, XeLaTeX with its default Latin Modern text font, reports no missing character.
\begin{equation}\label{tor:fun_orb2}
	u \mapsto \begin{cases} 
		\KK^\times, \text{ if } u \in \tau^\perp\\
		0 \text{ otherwise.} 
	\end{cases}
\end{equation}

\begin{theorem}\label{thm:mon}
Let $X = X_\sigma$ be an affine toric variety, the cone $\sigma$ contain a $k$-dimensional regular face $\tau \subset \sigma$, and $\{e_1^{(r)}, e_2^{(r)} \mid r = 1,\ldots,k\}$ be a set of Demazure roots compatible with $\tau$. Then the map $\KK[X] \rightarrow \KK[X] \otimes \KK[X]$ given by
\begin{equation}\label{thm:finalcom}
	\chi^u \mapsto \chi^u \otimes \chi^u \prod_{r = 1}^{k} (1 \otimes \chi^{e_1^{(r)}} + \chi^{e_2^{(r)}} \otimes 1)^{\langle p_r, u\rangle},
\end{equation}
where $\KK[X_\sigma] = \bigoplus_{u \in S_\sigma} \KK \chi^u$, defines a comultiplication that determines a monoid structure on~$X$ with group of invertible elements~$\GG$, where $\chi_r = \chi^{e_2^{(r)} - e_1^{(r)}}$. Moreover, the group of invertible elements coincides with the open subset $X_\tau$.
\end{theorem}

\begin{theorem}\label{idem_thm_1}
	In the notation of Theorem~\ref{thm:mon}, let $X_\sigma$ be a root monoid, $\tau$ the corresponding $k$-dimensional regular face of the cone $\sigma$ with primitive vectors $p_1, \ldots, p_k$ on the one-dimensional faces. Let $\gamma$ be a face of the cone $\sigma$. Denote by $E_\gamma$ the set of idempotents lying in the orbit $O_\gamma$. Then:
	\begin{enumerate}
		\item if $\tau \subseteq \gamma$, then $E_\gamma = \{x_\gamma \}$; 
		\item if there exists a ray $p_r \notin \gamma$ such that either corresponding Demazure roots $e_1^{(r)}, e_2^{(r)} \notin \gamma^\perp$, or $e_1^{(r)}, e_2^{(r)} \in \gamma^\perp$, then $E_\gamma = \varnothing$;
		\item if for every such ray $p_r \notin \gamma$ exactly one of the corresponding Demazure roots $e_1^{(r)}, e_2^{(r)}$ lies in $\gamma^\perp$, then \[E_\gamma = O_\gamma \cap \{\chi^u = 1\;\; \forall u \in \cone(\tau, \gamma)^\perp \cap S_\sigma\}.\]
	\end{enumerate}
\end{theorem}

\begin{proof}
	Note that $x \in X$ is an idempotent if and only if $\chi^u(x * x) = \chi^u(x)$ for every $u \in S_\sigma$.
	
	\bigskip
	
	(1) Suppose $\tau \subseteq \gamma$. Take $u \in \gamma^\perp \subseteq \tau^\perp$. Then $\langle p_r, u\rangle = 0$ for all $r = 1, \ldots, k$. By formula~\eqref{thm:finalcom} we have $\chi^u(x*x) = \chi^u(x)\cdot\chi^u(x)$. Hence $x$ can be an idempotent only if $\chi^u(x) = 1$ or $\chi^u(x) = 0$. Since $x \in O_\gamma, u \in \gamma^\perp$, by formula~\eqref{tor:fun_orb2} we know $\chi^u(x) \neq 0$, so $\chi^u(x) = 1$ and $x = x_\gamma$. Note that therefore there is at most one idempotent in the orbit $O_\gamma$. Let us check that $x_\gamma$ is indeed an idempotent. Now let $u \in \tau^\perp, u \notin \gamma^\perp$. Then $\chi^u(x_\gamma*x_\gamma) = \chi^u(x_\gamma)\cdot \chi^u(x_\gamma) = 0 \cdot 0 = 0 = \chi^u(x_\gamma)$. It remains to consider the case $u \notin \tau^\perp$. Recall that by Theorem~\ref{thm:mon}
	\begin{equation}\label{sum_form}
		\chi^u (x * y) = \sum \limits_{\ol{i} + \ol{j} = \langle \ol{p}, u \rangle } { \langle \ol{p}, u \rangle \choose \ol{i} } \chi^{u + \ol{i}\cdot \ol{e_2}}(x) \cdot  \chi^{u + \ol{j}\cdot \ol{e_1}} (y).
	\end{equation}
	
	Let us show that each summand in the right-hand side of~\eqref{sum_form} is equal to $0$. Recall that we denote $u + \ol{i}\cdot \ol{e_2} = u + i_1 e_2^{(1)} + \ldots + i_k e_2^{(k)}$. Then $\langle p_s, u + \ol{i}\cdot \ol{e_2} \rangle = \langle p_s, u\rangle - i_s$ for any $s = 1, \ldots, k$. Similarly $\langle p_s, u + \ol{j}\cdot \ol{e_1} \rangle = \langle p_s, u\rangle - j_s$. Since $\ol{i} + \ol{j} = \langle \ol{p}, u \rangle \neq \ol{0}$, at least one of $u + \ol{i}\cdot \ol{e_2}$ or $u + \ol{j}\cdot \ol{e_1}$ does not lie in $\tau^\perp$, and hence not in $\gamma^\perp$. Therefore $\chi^{u + \ol{i}\cdot \ol{e_2}}(x_\gamma) \cdot  \chi^{u + \ol{j}\cdot \ol{e_1}} (x_\gamma) = 0$. Then $\chi^u(x_\gamma*x_\gamma) = 0 = \chi^u(x_\gamma)$. Thus $x_\gamma$ is the unique idempotent in $O_\gamma$.
	
	\bigskip
	
	Let us divide case (2) into two subcases:\\
	(2.1) $p_r \notin \gamma$ and $e_1^{(r)}, e_2^{(r)} \notin \gamma^\perp$ for some $r$;\\
	(2.2) $p_r \notin \gamma$ and $e_1^{(r)}, e_2^{(r)} \in \gamma^\perp$ for some $r$, and there exists no $j$ such that $p_j \notin \gamma$ but ${e_1^{(j)}, e_2^{(j)} \notin \gamma^\perp}$.
	
	\medskip
	
	(2.1) Suppose $p_r \notin \gamma$ and $e_1^{(r)}, e_2^{(r)} \notin \gamma^\perp$. Then there exists a ray of the cone $\gamma$ with primitive vector $q$ such that $\langle q, e_1^{(r)} \rangle > 0$. Indeed, since $e_1^{(r)}\notin \gamma^\perp$, there exists $q \in \gamma$ with nonzero pairing. Since $e_1^{(r)}$ is a Demazure root, we have $\langle p_r, e_1^{(r)} \rangle = -1$, and with all other primitive vectors on rays of $\sigma$ the pairing is nonnegative, in particular $\langle q, e_1^{(r)} \rangle \geq 0$.  

	Next note that since $p_r$ is not a ray of $\gamma$, we have $\gamma^\perp \setminus p_r^\perp \neq \varnothing$. Take $u \in \gamma^\perp \setminus p_r^\perp$. Note that $i_r + j_r = \langle p_r, u \rangle \neq 0$, since $u \notin p_r^\perp$. Without loss of generality, assume $j_r \neq 0$. Then \[\langle q, u + \ol{j} \cdot \ol{e_1} \rangle = \langle q, u \rangle + j_1 \langle q, e_1^{(r)} \rangle + \ldots + j_k\langle q, e_1^{(r)} \rangle > 0.\] Indeed, $\langle q, u \rangle = 0$ since $q \in \gamma$, $u \in \gamma^\perp$, $j_r \langle q, e_1^{(r)} \rangle > 0$, and all other terms are nonnegative. Hence $u + \ol{j} \cdot \ol{e_1} \notin \gamma^\perp$. Thus for $x \in O_\gamma$, we get $\chi^{u + \ol{j} \cdot \ol{e_1}}(x)\cdot \chi^{u + \ol{i} \cdot \ol{e_2}}(x) = 0$. Then by~\eqref{sum_form}, $\chi^u(x*x) = 0$. On the other hand, $u \in \gamma^\perp$, whence $\chi^u(x) \neq 0$. Hence there are no idempotents in the orbit $O_\gamma$. 
	
	\bigskip
	
	For cases (2.2) and (3) we prove auxiliary lemmas. Recall that $\tau = \cone(p_1, \ldots, p_k)$; without loss of generality assume ${p_1, \ldots, p_s \notin \gamma}$, ${p_{s+1}, \ldots, p_k \in \gamma}$. 
	
	\begin{lemma}\label{lm:cone}
		Suppose that for each ray $p_r \notin \gamma$, where $r = 1, \ldots, s$, at least one of the corresponding Demazure roots $e_1^{(r)}$ or $e_2^{(r)}$ lies in $\gamma^\perp$. Then $\cone(p_1, \ldots, p_s, \gamma)$ is a face of the cone~$\sigma$.
	\end{lemma}
	
	\begin{proof}
	Denote by $e^{(r)}$ the Demazure root lying in $\gamma^\perp$, and by $\gamma_r$ the cone $\cone(p_1, \ldots, p_r, \gamma)$. We will prove by induction on $r$ that $\gamma_r$ is a face of $\sigma$. For $r = 0$ this is clear. Let $\gamma_r$ be a face of~$\sigma$, hence there exists $m_r \in M$ such that $\langle x, m_r\rangle = 0$ for all $x \in \gamma_r$, and $\langle x, m_r\rangle > 0$ for all other $y \in \sigma$. Take the Demazure root $e^{(r+1)} \in \gamma^\perp$ corresponding to the ray $p_{r+1}$. We will look for a similar vector for $\gamma_{r+1}$ of the form
	$$m_{r+1}:=m_r + \alpha e^{(r+1)}. $$ 
	Since the Demazure roots are compatible with the cone $\tau$, we have $\langle p_i,  e^{(r+1)} \rangle = -\delta_{i, r+1}$, where $\delta$ is the Kronecker symbol. Let $x \in \gamma_r$, that is, $x = x' + a_1 p_1 + \ldots + a_r p_r$, where ${x' \in \gamma}$. Since $e^{(r + 1)} \in \gamma^\perp$, we get $\langle x , m_{r+1} \rangle = \langle x , m_r \rangle + \langle x , \alpha e^{(r+1)} \rangle = 0$ for any $\alpha \in \ZZ$. For $x \notin \gamma_r$ we have $\langle x , m_{r+1} \rangle = \langle x , m_r \rangle  + \langle x , \alpha e^{(r+1)} \rangle > 0$, since the first term is strictly positive and the second is nonnegative. Next, note that $p_{r+1} \notin \gamma_r$, since $e^{(r+1)} \in \gamma_r^\perp$, but $e^{(r+1)} \notin p_{r+1}^\perp$. Therefore, $\langle p_{r+1} , m_r  \rangle > 0 $, and since $\langle p_{r+1} , e^{(r+1)} \rangle = -1$, we can choose $\alpha > 0$ such that $\langle p_{r+1} ,m_{r+1}\rangle = 0$. Now $\langle x, m_{r+1} \rangle = 0$ for all $x \in \cone(p_{r+1}, \gamma_r) = \gamma_{r+1}$, and $\langle x, m_{r+1} \rangle > 0$ for all other $x \in \sigma$. Hence $\gamma_{r+1}$ is also a face of the cone $\sigma$.
	\end{proof}

\begin{lemma}\label{lm:perpcone}
Suppose that for each of the rays $p_r \notin \gamma$, where $r = 1, \ldots, s$, at least one of the corresponding Demazure roots $e_1^{(r)}$ or $e_2^{(r)}$ lies in $\gamma^\perp$. Then for any $1 \leq r \leq s$ there exists $u \in \gamma^\perp \cap \sigma^\vee$ such that $\langle p_r, u \rangle = 1$, and $\langle p_j, u \rangle = 0$ for all $j \neq r,\; 1 \leq j \leq k$.
\end{lemma}

\begin{proof}
Let $e^{(r)} \in \gamma^\perp$ be the Demazure root corresponding to $p_r$. Note $\gamma^\perp \cap \{ \langle p_r, \cdot \rangle = 1\} \ni -e^{(r)}$. Take $w$ from the relative interior of the face $\gamma_s^\perp \cap \sigma^\vee$, where $\gamma_s = \cone(p_1, \ldots, p_s, \gamma) = \cone(\tau, \gamma)$ is a face of the cone $\sigma$ by Lemma~\ref{lm:cone}. Then $\langle p_j, w \rangle = 0$ and $\langle q, w \rangle > 0$ for the remaining primitive vectors $q$ on the rays of the cone $\sigma$ (not contained in $\gamma$). We will find $u$ in the form $u = -e^{(r)} + \alpha w$, where $\alpha \in \ZZ_{\geq 0}$. Note that $u \in \gamma^\perp$, and $\langle p_j, u \rangle = 0$ for $j \neq r$, $\langle p_r, u \rangle = 1$. Moreover, one can choose $\alpha$ in such a way that $\langle q, u \rangle \geq 0$ for the remaining primitive vectors $q$ on the rays of the cone $\sigma$, hence $u \in \sigma^\vee$.
\end{proof}

(2.2) Let $p_r \notin \gamma$, and $e_1^{(r)}, e_2^{(r)} \in \gamma^\perp$. Note that if there is no ray $p_j \notin \gamma$ for which both corresponding Demazure roots $e_1^{(j)}, e_2^{(j)} \notin \gamma^\perp$, then we may assume that for each of the rays $p_1, \ldots, p_s$, at least one of the corresponding Demazure roots $e_1^{(r)}$ or $e_2^{(r)}$ lies in $\gamma^\perp$. Therefore, we may apply Lemma~\ref{lm:perpcone}: there exists $u \in \gamma^\perp$ such that $\langle p_r, u \rangle = 1$, and $\langle p_j, u \rangle = 0$ for all $j \neq r$. Note that then $u + e_l^{(r)} \in \gamma^\perp \cap \tau^\perp$, $l = 1, 2$. Consider $x \in O_\gamma$, and suppose that $x$ is an idempotent. Let us prove that $\chi^{u + e_l^{(r)}}(x) = 1$. Indeed, by formula~\eqref{thm:finalcom} we have $\chi^{u + e_l^{(r)}}(x*x) = \chi^{u + e_l^{(r)}}(x)\cdot\chi^{u + e_l^{(r)}}(x)$. Therefore, $x$ can be an idempotent only if ${\chi^{u + e_l^{(r)}}(x) = 1}$ or $\chi^{u + e_l^{(r)}}(x) = 0$. Since $x \in O_\gamma,\; u + e_l^{(r)} \in \gamma^\perp$, it follows that $\chi^{u + e_l^{(r)}}(x) \neq 0$, hence ${\chi^{u + e_l^{(r)}}(x) = 1}$. On the other hand, since $\langle p_r, u \rangle = 1$, and $\langle p_j, u \rangle = 0$ for $j \neq r$, by formula~\eqref{thm:finalcom} we obtain $\chi^u(x*x) = \chi^u(x)\chi^{u + e_1^{(r)}}(x) + \chi^{u + e_2^{(r)}}\chi^u(x) = \chi^u(x) + \chi^u(x)$. That is, $\chi^u(x*x) = \chi^u(x)$ only when $\chi^u(x) = 0$, which contradicts $x \in O_\gamma$. Hence, $E_\gamma = \varnothing$.

\bigskip

(3) Now consider the case when for each $p_1, \ldots, p_s \notin \gamma$ exactly one of the Demazure roots $e_1^{(r)}$ or $e_2^{(r)}$ lies in $\gamma^\perp$. Note that in the expansion of $\chi^u(x*x)$ by formula~\eqref{thm:finalcom} the result does not change if we swap $e_1^{(r)}$ and $e_2^{(r)}$. Therefore, we may assume that $e_2^{(r)} \in \gamma^\perp$ and $e_1^{(r)} \notin \gamma^\perp$ for all $r = 1, \ldots, s$. Note that then for each root $e_1^{(r)}$, where $r = 1, \ldots, s$, there exists a ray of the cone~$\gamma$ and a primitive vector $q_r$ on this ray such that $\langle q_r, e_1^{(r)} \rangle \neq 0$. Since $p_r \notin \gamma$ while $q_r \in \gamma$, we have $p_r \neq q_r$, hence $\langle q_r, e_1^{(r)} \rangle > 0$. Moreover, since the set of Demazure roots is compatible with the cone $\tau$, i.e., $\langle p_i, e_1^{(r)} \rangle = -\delta_{ri}$ for any $i = 1, \ldots, k$, the vector $q_r$ does not coincide with any of $p_1, \ldots, p_k$. Therefore, the pairing of $q_r$ with all Demazure roots $\{e_1^{(r)}, e_2^{(r)} \mid r = 1,\ldots,k\}$ is nonnegative.

Consider $u \in \sigma^\vee \setminus \gamma^\perp$. We will check that $\chi^u(x*x) = 0$ for all $x \in O_\gamma$. Since $u \notin \gamma^\perp$, we obtain $\chi^u(x) = 0 = \chi^u(x*x)$ for all $x \in O_\gamma$. Denote by $\tau_0$ the face $\cone(p_{s+1}, \ldots, p_k) \subseteq \gamma$ of the cone $\tau$.

Let $u \in \sigma^\vee \setminus \gamma^\perp$ and $u \in \tau_0^\perp$. Choose the primitive vector $q$ on a ray of the cone $\gamma$ such that $\langle q, u \rangle > 0$. Note that then $q \notin \tau_0$, and hence $q$ has a nonnegative pairing with all Demazure roots $\{e_1^{(r)}, e_2^{(r)} \mid r = s+1,\ldots,k\}$. It follows that $\langle q + c_1 q_1 + \ldots + c_s q_s, u \rangle > 0$ for any choice of $c_r \in \ZZ_{\geq 0}$. Choose coefficients $c_r$ in such a way that $\langle q + c_1 q_1 + \ldots + c_s q_s, e_1^{(r)} \rangle > 0$ for all $r = 1, \ldots, s$; this is possible since $\langle q_r, e_1^{(r)} \rangle > 0$. Set $q' := q + c_1 q_1 + \ldots + c_s q_s \in \gamma$. Then $\langle q', u + j_1 e_1^{(1)} + \ldots + j_k e_1^{(k)} \rangle > 0$. Indeed, $\langle q', u \rangle > 0$, and the other terms are nonnegative. Thus, $u + j_1 e_1^{(1)} + \ldots + j_k e_1^{(k)} = u + \ol{j} \ol{e_1} \notin \gamma^\perp$, and by formula~\eqref{thm:finalcom} we obtain $\chi^u(x*x) = 0$ for all $x \in O_\gamma$.

Now let $u \in \sigma^\vee \setminus \gamma^\perp$ and $u \notin \tau_0^\perp$. Let us proceed with the reasoning analogous to that in case~(1). Recall that $\langle p_r, u + \ol{i}\cdot \ol{e_2} \rangle = \langle p_r, u \rangle - i_r = j_r$ and $\langle p_r, u + \ol{j}\cdot \ol{e_1} \rangle = \langle p_r, u \rangle - j_r = i_r$ for any $r = 1, \ldots, k$. Note that since $u \notin \tau_0^\perp$, it follows that $i_r + j_r = \langle p_r, u \rangle \neq 0$ for some $r \geq s+1$, hence $u + \ol{i}\cdot \ol{e_2}$ or $u + \ol{j}\cdot \ol{e_1}$ does not lie in $\tau_0^\perp$, and therefore not in $\gamma^\perp$, hence $\chi^{u + \ol{i}\cdot \ol{e_2}}(x_\gamma) \cdot \chi^{u + \ol{j}\cdot \ol{e_1}}(x_\gamma) = 0$. By formula~\eqref{thm:finalcom} we obtain $\chi^u(x*x) = 0$ for all $x \in O_\gamma$.

Now consider $u \in \gamma^\perp$. Then $\langle p_{s+1}, u \rangle = \ldots = \langle p_{k}, u \rangle = 0$. Then $j_{s+1} = \ldots = j_k = 0$ in formula~\eqref{thm:finalcom}. Let us show that $u + \ol{j} \ol{e_1} \notin \gamma^\perp$ for $\ol{j} \neq \ol{0}$. Suppose $j_r > 0$, where $r \leq s$. Then $\langle q_r, u + \ol{j} \ol{e_1} \rangle > 0$. Indeed, $j_r \langle q_r, e_1^{(r)} \rangle > 0$, while the other terms are nonnegative. Also note that $u + \ol{i} \ol{e_2} \in \gamma^\perp$. After expanding $\chi^u(x*x) = \chi^u(x)$ by formula~\eqref{thm:finalcom}, only one nonzero term remains on the left-hand side, corresponding to $\ol{j} = \ol{0}$. That is, we obtain the equality $\chi^u(x)\chi^{u + \langle \ol{p}, u \rangle \ol{e_2}}(x) = \chi^u(x)$. Since $u \in \gamma^\perp$, we have $\chi^u(x) \neq 0$, and we obtain $$\chi^{u + \langle \ol{p}, u \rangle \ol{e_2}}(x) = 1.$$ Denote $w = u + \langle \ol{p}, u \rangle \ol{e_2} = u + \langle p_1, u \rangle e_2^{(1)} + \ldots + \langle p_s, u \rangle e_2^{(s)}$. Notice that $w \in p_i^\perp$ for any $i = 1, \ldots, s$. Indeed, $\langle p_i, w \rangle = \langle p_i, u \rangle - \langle p_i, u \rangle = 0$. Thus, $w \in p_1^\perp \cap \ldots \cap p_s^\perp \cap \gamma^\perp$. In other words, $w \in \cone(\tau, \gamma)^\perp$. Moreover, choosing $u \in \cone(\tau, \gamma)^\perp$, we obtain as $w$ all possible elements of $\cone(\tau, \gamma)^\perp$, since in this case $w = u$.

Thus, the idempotents in $O_\gamma$ are the points $x$ satisfying $\chi^u(x) = 1$ for all $u \in \cone(\tau, \gamma)^\perp$.
\end{proof}

\end{document}
